% =====================================================================
% EXP-1 — Notas de la presentación (2 a 4 páginas sin contar anexos)
% Compilar: pdflatex notas -> biber notas -> pdflatex x2
% El guion de cada parte está en contenido/notas/: cada integrante edita
% SOLO los archivos de sus partes (ver partes.tex).
% =====================================================================
\documentclass[11pt,a4paper]{article}

\usepackage[T1]{fontenc}
\usepackage[utf8]{inputenc}
\usepackage[spanish,es-tabla,es-nodecimaldot,es-noshorthands]{babel}
\usepackage[default]{sourcesanspro}
\usepackage[a4paper,margin=2.1cm]{geometry}
\usepackage{microtype}
\usepackage{graphicx}
\usepackage{booktabs}
\usepackage{tabularx}
\usepackage{xltabular}
\usepackage{float}
\usepackage{csquotes}
\usepackage{bcm-comun}
\usepackage[backend=biber,style=ieee]{biblatex}
\addbibresource{referencias.bib}
\usepackage[hidelinks]{hyperref}

\setcounter{biburlnumpenalty}{100}
\setcounter{biburlucpenalty}{100}
\setcounter{biburllcpenalty}{100}

% Datos que comparten las diapositivas y las notas.
% El título es de trabajo: cámbienlo por el que decida el grupo.
\newcommand{\codigoactividad}{EXP-1}
\newcommand{\grupo}{BCM}
\newcommand{\tituloexposicion}{TCP, UDP y WebSockets}
\newcommand{\subtituloexposicion}{Exposición del primer corte}
\newcommand{\repositorio}{https://github.com/reyshadow0/isr701-exp-1}
\newcommand{\fechaexposicion}{lunes 23 de noviembre de 2026}
\newcommand{\docente}{Gleiston Cicerón Guerrero Ulloa, PhD}

% =====================================================================
% Reparto de la exposición: fuente única para diapositivas.tex y notas.tex
% =====================================================================
% Si cambian el orden, el título corto, quién expone o el tiempo de una
% parte, cámbienlo SOLO aquí: la agenda, la barra de tiempo de cada
% diapositiva y la tabla de las notas se actualizan solas.
% El total debe ser 12:00 (720 s); si pasa, LaTeX lo avisa en el log.

% \definirintegrante{id}{nombre corto}{nombre completo como en el SGA}
\definirintegrante{erwin}{Erwin}{Bueno Troya Erwin Daniel}
\definirintegrante{mario}{Mario}{Castro Coello Mario Sebastián}
\definirintegrante{cesar}{César}{Macías Vásquez César Abrahan}

% \definirparte{título corto}{id de quien expone}{segundos}
\definirparte{Apertura}{mario}{30}                   % 1
\definirparte{TCP y UDP}{erwin}{120}                 % 2
\definirparte{WebSockets}{mario}{120}                % 3
\definirparte{Ejemplo sencillo}{cesar}{120}          % 4
\definirparte{Panel de Erwin}{erwin}{90}             % 5
\definirparte{Panel de Mario}{mario}{90}             % 6
\definirparte{Panel de César}{cesar}{90}             % 7
\definirparte{Comparación de paneles}{erwin}{30}     % 8
\definirparte{Lecciones y cierre}{cesar}{30}         % 9

\bcmCalcularTotal


\hypersetup{pdftitle={\codigoactividad: notas de la presentación}, pdfauthor={Grupo \grupo}}

% Encabezado de cada parte del guion: \guion{n}
\newcommand{\guion}[1]{%
  \subsection*{\textcolor{bcmAcento}{0#1}\enspace\tituloParte{#1}}
  \vspace{-0.4em}
  {\small\color{bcmGris}Expone \textbf{\exponeParte{#1}}\enspace\textperiodcentered\enspace
    dura \duraParte{#1}\enspace\textperiodcentered\enspace empieza en \iniciaParte{#1}\par}
  \smallskip}

\begin{document}

% ---------------- Encabezado ----------------
\begin{tcolorbox}[enhanced, frame hidden, sharp corners, colback=bcmTinta, coltext=bcmPapel,
    left=12pt, right=12pt, top=10pt, bottom=10pt]
  {\small\color{bcmGris}Aplicaciones Distribuidas (ISR-701)\enspace\textperiodcentered\enspace UTEQ
    \enspace\textperiodcentered\enspace SPA 2026--2027\par}
  \smallskip
  {\Large\bfseries\codigoactividad: \tituloexposicion\par}
  {\color{bcmGris}Notas de la presentación\enspace\textperiodcentered\enspace Grupo \grupo
    \enspace\textperiodcentered\enspace\fechaexposicion\par}
  \medskip
  {\small\foreach \id in \bcmIntegrantes {\nombreCompleto{\id}\par}}
  \smallskip
  {\small\color{bcmGris}Docente: \docente\enspace\textperiodcentered\enspace
    Repositorio: \textcolor{bcmPapel}{\url{\repositorio}}\par}
\end{tcolorbox}

% ---------------- 1. Reparto ----------------
\section{Reparto y tiempos}

\renewcommand{\bcmFila}[1]{%
  \textcolor{bcmAcento}{\bfseries 0#1} & \tituloParte{#1} & \exponeParte{#1} &
  \iniciaParte{#1} & \duraParte{#1}\\}
\bcmPrepararFilas
\begin{table}[H]
  \centering
  \caption{Partes de la exposición, quién las expone y tiempos}
  \label{tab:reparto}
  \begin{tabular}{@{}rllrr@{}}
    \toprule
    & \textbf{Parte} & \textbf{Expone} & \textbf{Inicio} & \textbf{Duración}\\
    \midrule
    \the\bcmFilas
    \midrule
    & Preguntas & Todos & \bcmTiempo{\bcmTotalSeg} & 3:00\\
    \bottomrule
  \end{tabular}
\end{table}

Tiempo de cada integrante:
\foreach \id [count=\n] in \bcmIntegrantes {\ifnum\n>1 ; \fi\nombreCorto{\id}, \totalDe{\id}}.
La presentación dura \bcmTiempo{\bcmTotalSeg}; si pasa de 12:00 se interrumpe y lo que no se
alcanzó a presentar no se califica. Nadie lee de las diapositivas ni del celular, y en las
preguntas cualquiera de los tres responde sobre cualquier parte.

% ---------------- 2. Guion ----------------
\section{Guion de cada parte}

\guion{1}% ---------------------------------------------------------------------
% Parte 1 · Apertura · Mario · 0:30 (lo edita solo Mario)
% Escribe lo que vas a decir, en tus palabras y en 4 a 6 líneas: es tu
% guion, no una copia de las diapositivas. Cita cada dato (\cite[p.~X]{clave}).
% Preguntas guía (bórralas cuando termines):
%   - ¿Con qué pregunta o dato abres, y por qué le importa a quien escucha?
%   - ¿Con qué frase le pasas la palabra a Erwin?
% ---------------------------------------------------------------------
\pendiente{Mario}{guion de esta parte (4 a 6 líneas) con las citas de cada dato.}

\guion{2}% ---------------------------------------------------------------------
% Parte 2 · TCP y UDP · Erwin · 2:00 (lo edita solo Erwin)
% ---------------------------------------------------------------------
TCP y UDP están en la capa de transporte. Los dos mueven datos entre procesos de distintas máquinas, pero con una idea muy diferente.
TCP no envía nada hasta abrir la conexión con un handshake de tres pasos, y después numera los segmentos para que todo llegue completo y en el orden correcto \cite[p.~7]{barghash2022}.
Eso tiene un precio: cada envío espera su confirmación, y si algo se pierde hay que mandarlo otra vez, así que se demora más.
UDP no hace handshake ni revisa si los datagramas llegaron \cite[pp.~7--8]{barghash2022}. Es más liviano y por eso se usa donde perder un poco de información no importa tanto.
Pasa en una videollamada de Meet: si se pierde un pedazo de imagen ya no sirve pedirlo de nuevo, porque ese momento ya pasó.
Para el panel necesitamos que el servidor también pueda hablar primero. WebSocket hace eso y va sobre TCP, y es lo que les explica Mario.

\guion{3}% ---------------------------------------------------------------------
% Parte 3 · WebSockets · Mario · 2:00 (lo edita solo Mario)
% Escribe lo que vas a decir, en tus palabras y en 4 a 6 líneas: es tu
% guion, no una copia de las diapositivas. Cita cada dato (\cite[p.~X]{clave}).
% ---------------------------------------------------------------------
WebSocket se estandarizó en 2011. La idea es abrir una conexión entre el navegador y el servidor y no
cerrarla. Arranca con una petición HTTP que pide cambiar de protocolo, y si el servidor acepta, contesta
con un 101 \cite[secc.~2.4]{Murley2021}.

Con polling el navegador pregunta a cada rato. Según Murley y su equipo, cada ida y vuelta gasta cerca de
590 bytes solo en cabeceras \cite[p.~1197]{Murley2021}. Long polling reduce las preguntas, pero cada
mensaje sigue pagando ese costo.

En WebSocket cada mensaje lleva entre 2 y 14 bytes de cabecera \cite[secc.~2.4]{Murley2021}. En
tradingview una actualización pesa 82 bytes, y los autores calculan que con long polling serían unos 499
\cite[p.~1197]{Murley2021}.

\guion{4}En este ejemplo en Python probamos cómo funciona UDP en localhost usando dos hilos: uno para el servidor y otro para el cliente.
El servidor abre un socket de tipo SOCK\_DGRAM en el puerto 5002 y queda esperando con recvfrom, sin hacer handshake ni aceptar una conexión previa.
El cliente espera 0.3 segundos solo para asegurarse de que el servidor ya hizo bind y luego manda el paquete directo con sendto.
Cuando el servidor recibe el datagrama, lee el mensaje junto con la dirección de origen y devuelve la respuesta por la misma vía.
Con esto se ve lo principal de UDP: se ahorra el handshake y las retransmisiones, pero a cambio no garantiza que el paquete llegue ni que llegue rápido.
\guion{5}% ---------------------------------------------------------------------
% Parte 5 · Panel de Erwin · Erwin · 1:30 (lo edita solo Erwin)
% ---------------------------------------------------------------------
Mi panel está en la app de gestión de documentos de titulación, que es Spring Boot con Angular, y lo usa el estudiante mientras su director revisa el documento.
Cuando el director agrega una observación, lo devuelve o lo aprueba, el backend primero guarda el cambio y después manda un mensaje por WebSocket con lo que pasó y la hora.
Ese mensaje va a un canal que es solo de ese estudiante. Para conectarse necesita su token, y si intenta escuchar el canal de otro el servidor lo rechaza.
Antes había que recargar la página para enterarse. Ahora el aviso llega apenas pasa algo. Łasocha y Badurowicz compararon WebSocket con HTTP y midieron cuánto tarda cada uno en enviar datos \cite{lasocha2021}.
Si se cae la conexión, el indicador cambia de color y el cliente reintenta cada 5 segundos; al volver recarga el documento para no perder nada.
Si lo hiciera otra vez, guardaría los avisos en la base, así el estudiante vería también los que llegaron mientras estaba desconectado.

\guion{6}% ---------------------------------------------------------------------
% Parte 6 · Panel de Mario · Mario · 1:30 (lo edita solo Mario)
% Escribe lo que vas a decir, en tus palabras y en 4 a 6 líneas: es tu
% guion, no una copia de las diapositivas. Cita cada dato (\cite[p.~X]{clave}).
% ---------------------------------------------------------------------
Mi panel es del sistema de pre-sustentaciones y lo usa la coordinación para seguir las solicitudes.
Cuando un director cambia un estado desde su página, el servidor manda el evento por WebSocket y el panel
actualiza la fila sin recargar. Con polling tendría que preguntar a cada rato aunque nada cambie, y cada
pregunta gasta cerca de 590 bytes solo en cabeceras \cite[p.~1197]{Murley2021}.
Si se cae la conexión, el indicador pasa a ámbar y la tabla se queda con los últimos datos, que pueden
estar viejos. El panel reintenta a los 1, 2, 4 y 8 segundos y luego cada 10, y desde el cuarto intento se
pone rojo. Cuando vuelve el servidor, se reconecta solo y pide todas las solicitudes otra vez, así
recupera lo que cambió mientras estuvo caído.

\guion{7}En mi caso muestro el panel de monitoreo en tiempo real de SGIRI, pensado para que un operador supervise el estado de la red desde una sola pantalla.
Como se ve en la captura, el panel concentra cuatro métricas principales: latencia promedio, paquetes UDP con su porcentaje de pérdida, conexiones TCP activas y el estado general del servicio.
Debajo se actualiza cada segundo una gráfica de latencia junto con una tabla de eventos recientes que identifica el nodo, el protocolo y si hubo pérdida.
Aquí el WebSocket cumple el papel central: el servidor empuja las mediciones y alertas apenas ocurren, y el navegador las dibuja al vuelo sin tener que hacer peticiones HTTP repetitivas.
Si la conexión se cae, el indicador de estado avisa al operador, los eventos registran la pérdida y el cliente intenta reconectarse automáticamente para retomar la lectura apenas vuelve el enlace.
\guion{8}% ---------------------------------------------------------------------
% Parte 8 · Comparación de paneles · Erwin · 0:30 (lo edita solo Erwin)
% Escribe lo que vas a decir, en tus palabras y en 4 a 6 líneas: es tu
% guion, no una copia de las diapositivas. Cita cada dato (\cite[p.~X]{clave}).
% Preguntas guía (bórralas cuando termines):
%   - ¿En qué se parecen y en qué se diferencian los tres paneles?
%   - ¿Qué conclusión sale de ponerlos lado a lado?
% ---------------------------------------------------------------------
\pendiente{Erwin}{guion de esta parte (4 a 6 líneas) con las citas de cada dato.}

\guion{9}Para cerrar, nos quedamos con tres lecciones clave de toda la exposición.
Primero, ningún protocolo es mejor en abstracto: TCP prioriza entrega confiable, UDP evita el handshake y las retransmisiones, y WebSockets resuelve la comunicación bidireccional continua.
Segundo, toda decisión de transporte implica un compromiso: al quitar controles de red reducimos sobrecarga, pero la aplicación debe tolerar o gestionar posibles pérdidas.
Y tercero, en proyectos reales estos protocolos no compiten, se combinan según lo que necesita cada módulo.
En definitiva, elegir el protocolo correcto desde el diseño es lo que permite que un sistema distribuido funcione bien en producción.

% ---------------- 3. Preguntas ----------------
\section{Preguntas}

% ---------------------------------------------------------------------
% Preguntas (lo editan los tres)
% 1) Preguntas que esperan del docente o de otros equipos: cada una con
%    quién responde primero, aunque cualquiera debe poder hacerlo.
%    Practíquenlas cruzadas (Erwin responde de los paneles, César de TCP...).
% 2) La pregunta crítica que el grupo le hará a otro equipo (cuenta en D5).
% ---------------------------------------------------------------------
% xltabular: la tabla puede partirse entre páginas y repite el encabezado.
{\small\renewcommand{\arraystretch}{1.2}
\begin{xltabular}{\textwidth}{@{}>{\raggedright\arraybackslash}p{6.0cm}>{\raggedright\arraybackslash}p{2.4cm}X@{}}
  \caption{Preguntas previstas}\label{tab:preguntas}\\
  \toprule
  \textbf{Pregunta probable} & \textbf{Responde primero} & \textbf{Idea de la respuesta}\\
  \midrule
  \endfirsthead
  \toprule
  \textbf{Pregunta probable} & \textbf{Responde primero} & \textbf{Idea de la respuesta}\\
  \midrule
  \endhead
  \bottomrule
  \endlastfoot
  \pend{Grupo}{} & & \\
  \pend{Grupo}{} & & \\
  \pend{Grupo}{} & & \\
\end{xltabular}}

\noindent\textbf{Pregunta crítica para otro equipo:} \pend{Grupo}{a qué equipo y qué le van a preguntar}


% ---------------- 4. Aportes ----------------
\section{Tabla de aportes}

% Cada quien llena SU fila. Los aportes deben coincidir con los commits
% del repositorio (git log --author).
% xltabular: la tabla puede partirse entre páginas y repite el encabezado.
{\small\renewcommand{\arraystretch}{1.2}
\begin{xltabular}{\textwidth}{@{}>{\raggedright\arraybackslash}p{4.8cm}X>{\raggedright\arraybackslash}p{4.4cm}@{}}
  \caption{Aporte de cada integrante}\label{tab:aportes}\\
  \toprule
  \textbf{Integrante} & \textbf{Partes, figuras y evidencias} & \textbf{Referencias que aportó}\\
  \midrule
  \endfirsthead
  \toprule
  \textbf{Integrante} & \textbf{Partes, figuras y evidencias} & \textbf{Referencias que aportó}\\
  \midrule
  \endhead
  \bottomrule
  \endlastfoot
  \nombreCompleto{erwin} & \pend{Erwin}{} & \pend{Erwin}{}\\
  \nombreCompleto{mario} & \pend{Mario}{} & \pend{Mario}{}\\
  \nombreCompleto{cesar} & \pend{César}{} & \pend{César}{}\\
\end{xltabular}}

% ---------------- 5. Declaración de IA ----------------
\section{Declaración de uso de inteligencia artificial}

% Una fila por cada uso real, de cada integrante, escrita por el grupo.
% xltabular: la tabla puede partirse entre páginas y repite el encabezado.
{\small\renewcommand{\arraystretch}{1.2}
\begin{xltabular}{\textwidth}{@{}>{\raggedright\arraybackslash}p{2.8cm}>{\raggedright\arraybackslash}p{1.5cm}X>{\raggedright\arraybackslash}p{5.0cm}@{}}
  \caption{Herramientas de inteligencia artificial usadas como apoyo}\label{tab:declaracion-ia}\\
  \toprule
  \textbf{Herramienta} & \textbf{Versión} & \textbf{Para qué se usó} & \textbf{Cómo se verificó}\\
  \midrule
  \endfirsthead
  \toprule
  \textbf{Herramienta} & \textbf{Versión} & \textbf{Para qué se usó} & \textbf{Cómo se verificó}\\
  \midrule
  \endhead
  \bottomrule
  \endlastfoot
  Claude (Anthropic), en Claude Code & Opus 5.5 &
  Plantilla de las diapositivas (tema visual en Beamer) y de estas notas, estructura de carpetas,
  y propuesta de reparto de partes y tiempos. No redactó el contenido de las diapositivas ni del
  guion, ni buscó referencias. &
  \pend{Grupo}{cómo revisaron la plantilla y el reparto}\\
  \pend{Grupo}{} & & & \\
\end{xltabular}}

\printbibliography[heading=bibintoc, title={Referencias}]

% ---------------- Anexos ----------------
\clearpage
\appendix
\section{Esquema inicial a mano}

% Foto del primer esquema en papel, con los nombres de los tres y la fecha.
\begin{figure}[H]
  \centering
  \figura{evidencias/esquema-inicial.jpg}{0.8\textwidth}{9cm}{Foto del esquema inicial hecho a mano,
    con los nombres de los tres integrantes y la fecha.}
  \caption{Esquema inicial de la exposición, hecho a mano por el grupo}
  \label{fig:esquema}
\end{figure}

\section{Artículos citados, anotados}

% Al menos 3 capturas de artículos citados, subrayados o anotados, con la
% página visible. Pongan en el pie qué referencia es (\cite) y qué página.
\begin{figure}[H]
  \centering
  \figura{evidencias/articulo-1.png}{0.31\textwidth}{6cm}{Captura de un artículo citado, subrayado
    o anotado, con la página.}\hfill
  \figura{evidencias/articulo-2.png}{0.31\textwidth}{6cm}{Captura de un artículo citado, subrayado
    o anotado, con la página.}\hfill
  \figura{evidencias/articulo-3.png}{0.31\textwidth}{6cm}{Captura de un artículo citado, subrayado
    o anotado, con la página.}
  \caption{Artículos citados, anotados por el grupo: \pend{Grupo}{referencia y página de cada captura}}
  \label{fig:articulos}
\end{figure}

\end{document}
